\documentclass[11pt]{article}
\usepackage{amsmath,amssymb,amsthm}
\usepackage{booktabs}
\usepackage[hidelinks]{hyperref}

\emergencystretch=1.5em

\newtheorem{theorem}{Theorem}
\newtheorem{proposition}{Proposition}
\newtheorem{lemma}{Lemma}
\newtheorem{corollary}{Corollary}
\newtheorem{conjecture}{Conjecture}
\newtheorem{observation}{Observation}

\theoremstyle{remark}
\newtheorem*{remark}{Remark}

\title{Spectral Reformulations and Partial Structure Results for a Problem of\\ Green on Gowers Box Norms over Finite Fields}
\author{}
\date{}

\begin{document}
\maketitle
\vspace{-2em}

\begin{abstract}
We study an open problem of B.~Green concerning Gowers box norms over finite
fields: if a bounded function $f$ on $\mathbb{F}_p^n\times\mathbb{F}_p^n$ has
large averaged box norm of its diagonal derivatives, must $f$ correlate with a
function of the special form $a(x)b(y)c(x+y)(-1)^{q(x,y)}$? The problem is
open in both directions. We contribute five groups of results that do not
resolve it: an exact spectral reformulation of the hypothesis as a statement
about Fourier mass of the four-point box product on a zero-sum-frequency
subsystem; a multiplicative calculus showing that, inside the stated target
class, the hypothesis is precisely a Gowers $U^3$ condition on the
one-variable factor; a diagonal family showing that the case of maximal
hypothesis contains, as a special case, a twisted trilinear correlation
inequality for arbitrary unimodular data; an exact classification of one half
of a rigidity property of pointwise rank-one diagonal increments, together
with a sharpness example and a precise conjecture for the converse; and
small-field numerical experiments indicating a degrees-of-freedom-dominated
regime with no separation between structured and random diagonal data. No
disproof of the problem and no proof of its full correlation statement are
claimed.
\end{abstract}

\section{Introduction}\label{sec:intro}

Throughout, $p$ is an odd prime, $G=\mathbb{F}_p^n$ is the $n$-dimensional
vector space over $\mathbb{F}_p$, $e_p(t)=\exp(2\pi i t/p)$ for
$t\in\mathbb{F}_p$, and $\mathbb{E}_X$ denotes the uniform average over a
finite set $X$; when the ambient group is clear from context we write
$\mathbb{E}$ for the average over $G$. For $f:G\times G\to\mathbb{C}$ and
$h\in G$ the diagonal derivative is
\begin{equation}\label{eq:diagderiv}
\Delta_{(h,h)}f(x,y) \;=\; f(x+h,\,y+h)\,\overline{f(x,y)},
\end{equation}
and for $g:G\times G\to\mathbb{C}$ the (normalized) box norm of modulus four
is
\begin{equation}\label{eq:boxnorm}
\|g\|_{\square}^4 \;=\;
\mathbb{E}_{x_1,x_2,y_1,y_2}\,
g(x_1,y_1)\,\overline{g(x_2,y_1)}\,\overline{g(x_1,y_2)}\,g(x_2,y_2).
\end{equation}
For $w:G\to\mathbb{C}$ we write
\begin{equation}\label{eq:u2}
\|w\|_{U^2(G)}^4 \;=\;
\mathbb{E}_{m,r,s}\,
w(m)\,\overline{w(m+r)}\,\overline{w(m+s)}\,w(m+r+s)
\end{equation}
for the Gowers $U^2$ norm, where $\Delta_hc(t)=c(t+h)\overline{c(t)}$ denotes
the multiplicative derivative, and we set
$\|c\|_{U^3(G)}^8=\mathbb{E}_{h}\|\Delta_h c\|_{U^2(G)}^4$; the equivalence of
this expression with the usual cube average is proved in
Section~\ref{sec:calc}. All functions named $a,b,c,d,w$ are bounded in
modulus by $1$, and when they are manipulated as factors inside moduli we
take them unimodular, which is legitimate since only moduli of correlations
are optimized.

This paper is concerned with the following problem, posed by B.~Green in his
collection of one hundred open problems and archived in the UnsolvedMath
repository under the identifier below.

\begin{quote}\itshape
Let $p$ be an odd prime and suppose
$f:\mathbb{F}_p^n\times\mathbb{F}_p^n\to\mathbb{C}$ is bounded pointwise by
$1$. Suppose $\mathbb{E}_h\|\Delta_{(h,h)}f\|_\square^4\ge\delta$. Does $f$
correlate with a function of the form
$a(x)b(y)c(x+y)(-1)^{q(x,y)}$?
\end{quote}
\noindent
Here $a$, $b$, $c$ are one-variable functions on $G$ and $q$ is a polynomial;
the archived statement is \cite{unsolvedmath}. The phase $(-1)^{q(x,y)}$
admits two readings. Under the first, reading~(i), $(-1)^{q(x,y)}$ is
interpreted as a quadratic phase $e_p(Q(x,y))$ with
$Q:\mathbb{F}_p^{2n}\to\mathbb{F}_p$ a polynomial of degree at most $2$;
under the second, reading~(ii), $(-1)^{q(x,y)}$ is treated as a genuinely
$\{\pm1\}$-valued factor $\sigma$ of unspecified structure, with candidate
class
\begin{multline*}
C_{(ii)}=\bigl\{(x,y)\mapsto a(x)b(y)c(x+y)\sigma(x,y) \\
:\; |a|=|b|=|c|=1,\ \sigma:G\times G\to\{\pm1\}\bigr\}.
\end{multline*}
Every result below addresses reading~(i) only, through the class
\begin{multline}\label{eq:class}
\bigl\{\, (x,y)\mapsto a(x)\,b(y)\,c(x+y)\,e_p(Q(x,y)) \\
:\; a,b,c:G\to\mathbb{C},\ |a|=|b|=|c|=1,\ Q\ \text{quadratic} \,\bigr\},
\end{multline}
and nothing proved here extends to reading~(ii). The two candidate classes
are in fact incomparable once $n\ge2$: since $e_p(t)\in\{\pm1\}$ forces
$t\equiv0$ or $t\equiv(p-1)/2$ modulo $p$, a generic $\{\pm1\}$-valued phase
does not belong to \eqref{eq:class}; conversely $e_p(\lambda x_1y_2)$ with
$\lambda\neq0$ does not belong to $C_{(ii)}$, as the swap-asymmetry of its
mixed multiplicative difference shows. At $n=1$ every quadratic phase
absorbs into the three one-variable factors, so there $C=T_0$, where
$T_0:=\{a(x)b(y)c(x+y)\}$, and $T_0\subsetneq C_{(ii)}$.

The contributions are as follows.
\begin{enumerate}
\item \emph{Spectral reformulation}
(Proposition~\ref{prop:A}). The hypothesis
\[
\mathbb{E}_h\|\Delta_{(h,h)}f\|_\square^4\ge\delta
\]
is exactly the statement
that the four-point box product of $f$ carries Fourier $L^2$-mass at least
$\delta$ on the zero-sum-frequency subsystem of frequencies of $G^4$, a
subsystem of codimension $n$.
\item \emph{Multiplicative calculus of the target class}
(Lemma~\ref{lem:absorb}, Lemma~\ref{lem:noabsorb},
Proposition~\ref{prop:B}, Corollary~\ref{cor:reduction}). For $n=1$ every
quadratic phase absorbs into the three one-variable factors; for $n\ge2$ the
monomial $x_1y_2$ obstructs absorption, so the phase genuinely enlarges the
class. Inside the class, the box expression factors, and the hypothesis
becomes precisely $\|c\|_{U^3(G)}^8\ge\delta$; the inverse theorem for
$U^3$ then guarantees a quadratic-phase explanation for every class member
obeying the hypothesis.
\item \emph{The diagonal family} (Observation~\ref{obs:C}). Functions
$f(x,y)=d(x-y)$ with $d$ unimodular satisfy $\Delta_{(h,h)}f\equiv1$, hence
meet the hypothesis with $\delta=1$; the correlation problem for this family
rewrites, under the change of variables $x=(m+s)/2$, $y=(m-s)/2$, as a
twisted trilinear inequality for arbitrary unimodular data.
\item \emph{Rigidity} (Theorem~\ref{thm:D}, Conjecture~\ref{conj:Dp}). Every
member of the extended class has pointwise rank-one diagonal increments,
with an explicit splitting; a cubic mixed perturbation destroys this
structure for most directions. The exact converse classification is stated
as an open conjecture, the defining rank-one increment law being the only
currently verified invariant.
\item \emph{Numerics} (Section~\ref{sec:numerics}). Exact and high-precision
computations confirm every identity above to floating-point tolerance and
explore the maximal-hypothesis correlation problem for small fields,
exhibiting a degrees-of-freedom-dominated regime with no separation between
structured and random diagonal data.
\end{enumerate}

An honest status declaration is in order. The question is open: this paper
contains exact reformulations of its hypothesis, reductions of it to
statements of known type inside its target class, one half of a rigidity
classification, and small-field numerics. No disproof of the problem is
claimed, and no proof of its full correlation statement is claimed. The
precise gap can be stated in one sentence: the Fourier-mass condition
\eqref{eq:spectral} on $B_f$ is \emph{not known to imply} algebraic
correlation with members of the infinite-dimensional target class, and no
implication from the spectral condition to correlation is currently known.
These sentences — like every statement below of the form ``remains open,''
``no known implication,'' ``no counterexample is known,'' and similar — are
literature-status assessments rather than deductions from the preceding
results.
All numerical residuals quoted in this paper were independently confirmed to
floating-point tolerance at most $10^{-15}$ by an external deterministic
verification suite, \texttt{verify\_solution.py}.

\section{A spectral reformulation}\label{sec:spectral}

Identify the dual group of $G$ with $G$ itself, coordinatewise, and write
characters in the form $z\mapsto e_p(\xi\cdot z)$. For
$f:G\times G\to\mathbb{C}$ define the four-point box product
\begin{equation}\label{eq:boxproduct}
B_f(x_1,x_2,y_1,y_2)\;=\;
f(x_1,y_1)\,\overline{f(x_2,y_1)}\,\overline{f(x_1,y_2)}\,f(x_2,y_2),
\end{equation}
a function on $G^4$, and its Fourier coefficients
\begin{equation}\label{eq:fouriercoef}
\widehat{B_f}(\xi)\;=\;\mathbb{E}_{z\in G^4}\,
B_f(z)\,\overline{e_p(\xi\cdot z)},
\qquad \xi=(\xi_1,\xi_2,\xi_3,\xi_4)\in \widehat{G^4}.
\end{equation}
Let
\begin{equation}\label{eq:Xi0}
\Xi_0 \;=\; \bigl\{\xi=(\xi_1,\xi_2,\xi_3,\xi_4)\in\widehat{G^4} :
\xi_1+\xi_2+\xi_3+\xi_4=0 \text{ in } \widehat{G}\bigr\},
\end{equation}
the zero-sum-frequency subsystem, of codimension $n$ in $\widehat{G^4}$.

\begin{proposition}\label{prop:A}
For every $f:G\times G\to\mathbb{C}$,
\begin{equation}\label{eq:spectral}
\mathbb{E}_h\,\bigl\|\Delta_{(h,h)}f\bigr\|_\square^4
\;=\;
\sum_{\xi\in\Xi_0}\bigl|\widehat{B_f}(\xi)\bigr|^2 .
\end{equation}
\end{proposition}

\begin{proof}
Fix $h\in G$ and write $g_h=\Delta_{(h,h)}f$, so that
$g_h(x,y)=f(x+h,y+h)\overline{f(x,y)}$. Insert this into the box expression
\eqref{eq:boxnorm} and expand each of the four factors:
\begin{align*}
g_h(x_1,y_1)&=f(x_1+h,y_1+h)\,\overline{f(x_1,y_1)},\\
\overline{g_h(x_2,y_1)}&=\overline{f(x_2+h,y_1+h)}\,f(x_2,y_1),\\
\overline{g_h(x_1,y_2)}&=\overline{f(x_1+h,y_2+h)}\,f(x_1,y_2),\\
g_h(x_2,y_2)&=f(x_2+h,y_2+h)\,\overline{f(x_2,y_2)}.
\end{align*}
Multiplying the four lines and comparing with \eqref{eq:boxproduct} term by
term, the four unconjugated shifted factors assemble to
$B_f(x_1+h,x_2+h,y_1+h,y_2+h)$ and the four conjugated unshifted factors
assemble to $\overline{B_f(x_1,x_2,y_1,y_2)}$. Hence, with
$z=(x_1,x_2,y_1,y_2)$ and $u=(h,h,h,h)$,
\begin{equation}\label{eq:shiftstep}
\bigl\|\Delta_{(h,h)}f\bigr\|_\square^4
\;=\;\mathbb{E}_{z\in G^4}\,
B_f(z+u)\,\overline{B_f(z)}.
\end{equation}
Expanding $B_f$ in its Fourier series
$B_f=\sum_{\xi}e_p(\xi\cdot z)\,\widehat{B_f}(\xi)$, whose characters are
orthonormal for the normalized inner product,
$\mathbb{E}_z\,e_p(\xi\cdot z)\overline{e_p(\eta\cdot z)}
=\mathbf{1}_{\xi=\eta}$, the right side of \eqref{eq:shiftstep} equals
\[
\sum_{\xi,\eta\in\widehat{G^4}}
\widehat{B_f}(\xi)\,\overline{\widehat{B_f}(\eta)}\,
e_p(\xi\cdot u)\,\mathbb{E}_z e_p((\xi-\eta)\cdot z)
\;=\;
\sum_{\xi\in\widehat{G^4}}\bigl|\widehat{B_f}(\xi)\bigr|^2\,e_p(\xi\cdot u).
\]
Finally, $\xi\cdot u=h\cdot(\xi_1+\xi_2+\xi_3+\xi_4)$, so averaging over
$h\in G$ and using orthogonality
$\mathbb{E}_h\,e_p(h\cdot\kappa)=\mathbf{1}_{\kappa=0}$ retains precisely the
frequencies with $\xi_1+\xi_2+\xi_3+\xi_4=0$, that is, the sum appearing in
\eqref{eq:spectral}.
\end{proof}

\begin{remark}\label{rem:spectral}
The identity \eqref{eq:spectral} exhibits the left side as a nonnegative
spectral mass. Consequently the hypothesis of the problem says exactly that
the four-point box product $B_f$ carries Fourier $L^2$-mass at least
$\delta$ on the zero-sum-frequency subsystem $\Xi_0$, a subsystem of
codimension $n$ in $\widehat{G^4}$; no mass on complementary frequencies
contributes to the hypothesis. The identity was confirmed to floating-point
tolerance $10^{-15}$ at the three group sizes listed in
Table~\ref{tab:ident}, together with the Parseval diagnostic that the total
Fourier mass of $B_f$ equals $1$ at all three sizes.
\end{remark}

\begin{table}[ht]
\centering
\small
\begin{tabular}{lccc}
\toprule
$(p,n)$ &
$\mathbb{E}_h\|\Delta_{(h,h)}f\|_\square^4$ &
$\sum_{\xi\in\Xi_0}|\widehat{B_f}(\xi)|^2$ &
total mass \\
\midrule
$(3,1)$ & $0.6049382716049383$ & $0.6049382716049382$ & $1.0000000$ \\
$(3,2)$ & $0.307422648986435$  & $0.3074226489864356$ & $1.0000000$ \\
$(5,1)$ & $0.45728$            & $0.4572800000000002$ & $1.0000000$ \\
\bottomrule
\end{tabular}
\caption{Identity checks for Proposition~\ref{prop:A} on random
$\{\pm1\}$-valued data, with the Parseval diagnostic.}
\label{tab:ident}
\end{table}

\section{Multiplicative calculus of the target class}\label{sec:calc}

Two lemmas determine to what extent the quadratic phase in the class
\eqref{eq:class} is redundant.

\begin{lemma}[Absorption in one variable]\label{lem:absorb}
Let $n=1$ and let $p$ be odd. Every polynomial $Q(x,y)$ of degree at most
$2$ on $\mathbb{F}_p\times\mathbb{F}_p$ can be written as
\[
Q(x,y)\;=\;\alpha(x)+\beta(y)+\gamma(x+y)+\kappa
\]
with $\alpha,\beta,\gamma$ polynomials of degree at most $2$ in one variable
and $\kappa$ a constant.
\end{lemma}

\begin{proof}
Collect coefficients and write, with all arithmetic in $\mathbb{F}_p$,
\[
Q(x,y)\;=\;A x^2+B xy+C y^2+Dx+Ey+F.
\]
The identity $xy=\bigl((x+y)^2-x^2-y^2\bigr)/2$, legitimate since $2$ is
invertible modulo $p$, converts the mixed monomial into separated ones.
Setting $\lambda=B/2$ and
\begin{gather*}
\alpha(x)=\Bigl(A-\frac{B}{2}\Bigr)x^2+Dx,\qquad
\beta(y)=\Bigl(C-\frac{B}{2}\Bigr)y^2+Ey,\\
\gamma(s)=\frac{B}{2}\,s^2,\qquad
\kappa=F,
\end{gather*}
direct expansion gives
$\alpha(x)+\beta(y)+\gamma(x+y)+\kappa
=\bigl(A-\tfrac{B}{2}\bigr)x^2+\bigl(C-\tfrac{B}{2}\bigr)y^2
+\tfrac{B}{2}(x+y)^2+Dx+Ey+F
=Ax^2+Bxy+Cy^2+Dx+Ey+F=Q(x,y)$,
because $(x+y)^2=x^2+2xy+y^2$ contributes exactly
$\tfrac{B}{2}\cdot 2xy=Bxy$. Each of $\alpha,\beta,\gamma$ has degree at
most $2$.
\end{proof}

\begin{lemma}[Obstruction to absorption in several variables]\label{lem:noabsorb}
Let $n\ge2$. The monomial $m(x,y)=x_1y_2$ on $G\times G$ cannot be written
as
$m(x,y)=\alpha(x)+\beta(y)+\gamma(x+y)+\kappa$
with $\alpha,\beta,\gamma:G\to\mathbb{F}_p$ arbitrary functions and
$\kappa$ constant. Consequently, for $n\ge2$ there exist quadratic phases
$e_p(Q)$ that are not products of a unimodular function of $x$, a function
of $y$, and a function of $x+y$; the phase factor genuinely enlarges the
target class.
\end{lemma}

\begin{proof}
Suppose such a representation existed, with all arithmetic in
$\mathbb{F}_p$. For a fixed increment $w\in G$ define the difference
$G_w(s)=\gamma(s+w)-\gamma(s)$. Apply the mixed finite difference, with
increment $w$ in the $x$-variable and $z$ in the $y$-variable, to the
purported identity. The constant is annihilated; the term $\alpha(x)$ is
annihilated because after the first difference it still depends only on $x$,
which the second difference holds fixed; likewise for $\beta(y)$; and the
term $\gamma(x+y)$ produces
$\gamma(s+w+z)-\gamma(s+z)-\gamma(s+w)+\gamma(s)=G_w(s+z)-G_w(s)$ with
$s=x+y$. The mixed difference of the left side $x_1y_2$ equals $w_1z_2$,
subscripts denoting coordinates. Hence
\begin{equation}\label{eq:diffeq}
G_w(s+z)-G_w(s)\;=\;w_1z_2
\qquad\text{for all } s,z\in G.
\end{equation}
Setting $s=0$ in \eqref{eq:diffeq} gives $G_w(z)=G_w(0)+w_1z_2$ for every
$z$; since $G_w(0)=\gamma(w)-\gamma(0)$,
\begin{equation}\label{eq:gammarel}
\gamma(s+w)-\gamma(s)\;=\;\gamma(w)-\gamma(0)+w_1s_2
\qquad\text{for all } s,w\in G.
\end{equation}
Interchanging the roles of $s$ and $w$ in \eqref{eq:gammarel} and comparing
the two expressions for $\gamma(s+w)$ yields
$w_1s_2=s_1w_2$ for all $s,w\in G$. Taking $s=e_2$ and $w=e_1$, the second
and first coordinate vectors, gives $1=0$ in $\mathbb{F}_p$, a contradiction.

An algebraic rank computation over $\mathbb{F}_5$ confirms the obstruction
quantitatively: the linear system expressing a prescribed quadratic form as
$\alpha(x)+\beta(y)+\gamma(x+y)+\kappa$ is consistent exactly when the
cross-coordinate block of the prescribed form is compatible with the blocks
realizable by $\alpha$, $\beta$, and $\gamma$; forcing the incompatible
requirement in which the prescribed cross-coordinate block is
$\left[\begin{smallmatrix}0&1\\1&0\end{smallmatrix}\right]$ while the
realizable system supplies
$\left[\begin{smallmatrix}1&0\\0&0\end{smallmatrix}\right]$
produces an inconsistent linear system whose rank falls one short of the
required value.
\end{proof}

The next proposition records that, against the box pattern, the factors
$a(x)b(y)$ cancel identically.

\begin{proposition}[Factorization]\label{prop:B}
If $a,b:G\to\mathbb{C}$ are unimodular and $w:G\to\mathbb{C}$ is arbitrary,
then
\[
\bigl\|\,a(x)\,b(y)\,w(x+y)\bigr\|_\square^4\;=\;\|w\|_{U^2(G)}^4 .
\]
\end{proposition}

\begin{proof}
Write $g(x,y)=a(x)b(y)w(x+y)$ and expand \eqref{eq:boxnorm}. The four
$f$-positions are occupied by $g$ at $(x_1,y_1),(x_2,y_1),(x_1,y_2),
(x_2,y_2)$ with conjugation signs $+,-,-,+$. Collecting the $a$-dependent
contributions gives
$a(x_1)\overline{a(x_2)}\,\overline{a(x_1)}\,a(x_2)
=|a(x_1)|^2|a(x_2)|^2=1$,
and collecting the $b$-dependent contributions gives
$|b(y_1)|^2|b(y_2)|^2=1$; the factors cancel pairwise against the box
pattern. What remains is
$w(s_{11})\overline{w(s_{21})}\,\overline{w(s_{12})}\,w(s_{22})$ with
$s_{ij}=x_i+y_j$. Now perform the change of variables
$x_2=x_1+r$, $y_2=y_1+s$, a bijection of $G^4$, and set $m=x_1+y_1$; the
four sums become $(m,\ m+r,\ m+s,\ m+r+s)$, so they sweep every additive
quadruple counted by \eqref{eq:u2} exactly once. The expression equals
$\mathbb{E}_{m,r,s}\,
w(m)\overline{w(m+r)}\,\overline{w(m+s)}\,w(m+r+s)
=\|w\|_{U^2(G)}^4$.
\end{proof}

\begin{corollary}[Reduction inside the class]\label{cor:reduction}
Let $f(x,y)=a(x)b(y)c(x+y)$ with $a,b,c:G\to\mathbb{C}$ unimodular. Then
\begin{equation}\label{eq:u3reduction}
\mathbb{E}_h\,\bigl\|\Delta_{(h,h)}f\bigr\|_\square^4
\;=\;\mathbb{E}_h\,\|\Delta_{2h}c\|_{U^2(G)}^4
\;=\;\|c\|_{U^3(G)}^8 .
\end{equation}
Thus, restricted to the target class \eqref{eq:class} without phase, the
hypothesis of the problem is precisely a $U^3$ condition on the one-variable
factor $c$.
\end{corollary}

\begin{proof}
Unrolling \eqref{eq:diagderiv},
\[
\Delta_{(h,h)}f(x,y)
=a(x+h)\overline{a(x)}\cdot b(y+h)\overline{b(y)}
\cdot c(x+y+2h)\overline{c(x+y)} .
\]
Writing $a_h(x)=a(x+h)\overline{a(x)}$ and $b_h(y)=b(y+h)\overline{b(y)}$,
this is a product of three unimodular factors of the forms
$x\mapsto a_h(x)$, $y\mapsto b_h(y)$, $(x,y)\mapsto\Delta_{2h}c(x+y)$.
Proposition~\ref{prop:B} applied with $w=\Delta_{2h}c$ gives
$\|\Delta_{(h,h)}f\|_\square^4=\|\Delta_{2h}c\|_{U^2(G)}^4$. Since
$h\mapsto 2h$ is a bijection of $G$ for odd $p$, averaging over $h$ may be
reparametrized, giving the first equality of \eqref{eq:u3reduction}.

For the second equality, expand the right side:
\begin{multline*}
\mathbb{E}_{h,m,r,s}\,
c(m+h)\overline{c(m)}\,
\overline{c(m+r+h)}\,c(m+r)\,
\overline{c(m+s+h)}\,c(m+s)\\
\times\; c(m+r+s+h)\,\overline{c(m+r+s)} .
\end{multline*}
Relabeling $x=m$, $h_1=h$, $h_2=r$, $h_3=s$ and reordering the eight
factors produces exactly the cube average
\begin{multline*}
\mathbb{E}_{x,h_1,h_2,h_3}\,
c(x+h_1+h_2+h_3)\,\overline{c(x+h_1+h_2)}\,
\overline{c(x+h_1+h_3)}\,c(x+h_1)\\
\times\;\overline{c(x+h_2+h_3)}\,c(x+h_2)\,c(x+h_3)\,\overline{c(x)},
\end{multline*}
which is $\|c\|_{U^3(G)}^8$.
\end{proof}

The residual between the two sides of \eqref{eq:u3reduction} was measured to
be $2.220\times10^{-16}$ at $(p,n)=(3,1)$. The reduction combines with the
inverse theorem for the Gowers $U^3$ norm over finite-field vector spaces
\cite{greentao2010}: every member of the target class obeying the hypothesis
with parameter $\delta$ has its one-variable factor correlate with a
function carrying a quadratic phase twist. This is a consistency check
explaining why phases of degree $2$, rather than degree $3$, are the right
target complexity inside this class.

\section{The diagonal family}\label{sec:diagonal}

\begin{observation}\label{obs:C}
Let $d:G\to\mathbb{C}$ be unimodular, $|d|\equiv1$; in particular $d$ is
bounded by $1$. Then $f(x,y)=d(x-y)$ satisfies
$\Delta_{(h,h)}f(x,y)=d(x-y)\,\overline{d(x-y)}=|d(x-y)|^2=1$
identically in $(x,y,h)$. Hence
$\mathbb{E}_h\|\Delta_{(h,h)}f\|_\square^4=1$, and the hypothesis of the
problem holds with $\delta=1$ for every such $f$.
\end{observation}

Consequently, as transmitted, the problem contains at $\delta=1$ the
following special case: every bounded one-variable function $d$ must
correlate with the target class \eqref{eq:class}. Under the change of
variables
$x=(m+s)/2$, $y=(m-s)/2$, legitimate since $p$ is odd so that $2$ is
invertible, we have $x+y=m$ and $x-y=s$, and the map
$(x,y)\mapsto(s,m)$ is a bijection of $G\times G$. The correlation of
$f=d(x-y)$ against a member of \eqref{eq:class} becomes
\begin{equation}\label{eq:twistedcorr}
\Bigl|\;
\mathbb{E}_{s,m\in G}\;
d(s)\,
\overline{a\bigl(\tfrac{m+s}{2}\bigr)}\,
\overline{b\bigl(\tfrac{m-s}{2}\bigr)}\,
\overline{c(m)}\,
e_p\!\bigl(-Q\bigl(\tfrac{m+s}{2},\tfrac{m-s}{2}\bigr)\bigr)
\;\Bigr|.
\end{equation}

Two qualifications govern how much freedom the twist
$Q\bigl(\tfrac{m+s}{2},\tfrac{m-s}{2}\bigr)$ contributes, according to $n$.
When $n\ge2$, the pullback of a quadratic form under an invertible linear
change of variables is again a quadratic form, and conversely every
quadratic form in $(s,m)$ arises this way; since
Lemma~\ref{lem:noabsorb} shows that cross-coordinate quadratic phases such
as $e_p(\lambda x_1y_2)$ do not absorb into $a,b,c$, the twist in
\eqref{eq:twistedcorr} ranges over all quadratic forms in $(s,m)$, including
forms with nonzero mixed $sm$-terms. When $n=1$, every quadratic form is of
the shape analyzed in Lemma~\ref{lem:absorb}; although the pulled-back form
$A\bigl(\tfrac{m+s}{2}\bigr)^2+B\bigl(\tfrac{m+s}{2}\bigr)
\bigl(\tfrac{m-s}{2}\bigr)+C\bigl(\tfrac{m-s}{2}\bigr)^2
=\bigl[(A{+}B{+}C)m^2+2(A{-}C)ms+(A{-}B{+}C)s^2\bigr]/4$
may display a mixed $ms$-coefficient when $A\ne C$, the entire phase absorbs
into the three unimodular factors by Lemma~\ref{lem:absorb}, so modulo this
absorption no residual quadratic twist remains and \eqref{eq:twistedcorr}
reduces to a purely trilinear correlation.

This yields the following equivalent formulation of the case $\delta=1$.

\emph{Twisted trilinear inequality.} For every unimodular $d:G\to\mathbb{C}$
there exist unimodular $a,b,c$ (and, for $n\ge2$, possibly a quadratic twist
as in \eqref{eq:twistedcorr}) such that
\[
\bigl|\mathbb{E}_{s,m\in G}\;
d(s)\,\overline{a((m+s)/2)}\,\overline{b((m-s)/2)}\,\overline{c(m)}\,
e_p(-Q((m+s)/2,(m-s)/2))\bigr|
\]
is bounded below by a positive quantity independent of $d$.

\noindent
Whether such a uniform lower bound holds is open; it is recorded as
Open Problem~O1 in Section~\ref{sec:open}.

\section{Rigidity of pointwise rank-one diagonal increments}\label{sec:rigidity}

Say that $f:G\times G\to\mathbb{C}$ has \emph{pointwise rank-one diagonal
increments} if for every $h\in G$ there exist functions
$\alpha_h,\beta_h:G\to\mathbb{C}$ with
$\Delta_{(h,h)}f(x,y)=\alpha_h(x)\,\beta_h(y)$ for all $(x,y)$.

\begin{theorem}[Sufficiency]\label{thm:D}
Let $a,b,d:G\to\mathbb{C}$ be unimodular, let $Q$ on $G\times G$ be a
polynomial of degree at most $2$, written in block form as
$Q(z)=\tfrac12 z^{T}Mz+l^{T}z+\kappa$ with
$z=(x,y)\in G\times G$,
$M=\left[\begin{smallmatrix}M_{xx}&M_{xy}\\ M_{yx}&M_{yy}\end{smallmatrix}\right]$
symmetric, $l\in(G\times G)^{*}\cong\mathbb{F}_p^{2n}$, and $\kappa$
constant. Then
$f(x,y)=a(x)b(y)d(x-y)e_p(Q(x,y))$ has pointwise rank-one diagonal
increments; explicitly, for each $h\in G$,
\begin{equation}\label{eq:rankone}
\Delta_{(h,h)}f(x,y)\;=\;\alpha_h(x)\,\beta_h(y)\;e_p(\omega_h),
\end{equation}
where
$\alpha_h(x)=a(x+h)\overline{a(x)}\;e_p\bigl((M_{xx}+M_{xy})h\bigr)^{T}x$,
$\beta_h(y)=b(y+h)\overline{b(y)}\;e_p\bigl((M_{yx}+M_{yy})h\bigr)^{T}y$,
and $\omega_h=\tfrac12(h,h)^{T}M(h,h)+l^{T}(h,h)$.
\end{theorem}

\begin{proof}
For $\delta=(h,h)$, the increment of the quadratic part is
\[
Q(z+\delta)-Q(z)
=\tfrac12\bigl(\delta^{T}Mz+z^{T}M\delta+\delta^{T}M\delta\bigr)+l^{T}\delta
=(M\delta)^{T}z+\tfrac12\delta^{T}M\delta+l^{T}\delta ,
\]
using symmetry of $M$. Since $M\delta=
\bigl((M_{xx}+M_{xy})h,\ (M_{yx}+M_{yy})h\bigr)$, writing
$P=M_{xx}+M_{xy}$, $R=M_{yx}+M_{yy}$, and
$\omega_h=\tfrac12(h,h)^{T}M(h,h)+l^{T}(h,h)$, we obtain the exact identity
\[
Q(x+h,y+h)-Q(x,y)\;=\;(Ph)^{T}x+(Rh)^{T}y+\omega_h ,
\]
whose mixed $x{\cdot}y$ coefficient vanishes identically: the increment is
affine separately in $x$ and in $y$. Combining this with the triviality
$d(x+h-(y+h))\,\overline{d(x-y)}
=d(x-y)\,\overline{d(x-y)}=1$ gives
\[
\Delta_{(h,h)}f(x,y)
=a(x+h)\overline{a(x)}\,e_p\bigl((Ph)^{T}x\bigr)\;\cdot\;
b(y+h)\overline{b(y)}\,e_p\bigl((Rh)^{T}y\bigr)\;\cdot\;e_p(\omega_h),
\]
which is \eqref{eq:rankone}.
\end{proof}

At $(p,n)=(5,1)$ the splitting \eqref{eq:rankone} was verified over all
triples $(h,x,y)$ with maximal deviation $5.118\times10^{-16}$ from the
pointwise product of its two factors. The structure is also sharp within the
extended class.

\begin{corollary}[Sharpness]\label{cor:sharp}
Multiply the function of Theorem~\ref{thm:D} at $(p,n)=(5,1)$ by the cubic
mixed phase $e_p(xy^{2})$. Then for every $h\ne0$ the increment acquires the
mixed term $2h\,xy$: indeed
\[
(x+h)(y+h)^2-xy^2\;=\;2h\,xy+h\,y^{2}+2h^{2}y+h^{2}x+h^{3},
\]
and $2h\not\equiv0$ modulo the odd prime $p$, so the increment cannot split
as a function of $x$ times a function of $y$. Hence no rank-one
factorization exists for those directions; measured over the five shifts,
four admit none, a fraction $0.8$ of all directions.
\end{corollary}

The twist of Corollary~\ref{cor:sharp} does more than destroy rank-one
factorization: it destroys the \emph{hypothesis} itself.

\begin{proposition}[Collapse under a cubic mixed twist]\label{prop:collapse}
Let $g(x,y)=a(x)b(y)d(x-y)e_p(Q)$ lie in the family of
Theorem~\ref{thm:D}, and let $f=g\,e_p(x_0y_0^{2})$, where $x_0,y_0$ are
the zeroth coordinates. Then
\begin{equation}\label{eq:collapse}
\bigl\|\Delta_{(h,h)}f\bigr\|_\square^4\;=\;
\begin{cases}
1, & h_0=0,\\[2pt]
p^{-1}, & h_0\neq0,
\end{cases}
\qquad\text{whence}\qquad
\mathbb{E}_h\,\bigl\|\Delta_{(h,h)}f\bigr\|_\square^4\;=\;\frac{2p-1}{p^{2}},
\end{equation}
independently of $n$ and of the choice of $a,b,d,Q$; and $f$ violates the
hypothesis of the problem for every $\delta>(2p-1)/p^{2}$.
\end{proposition}

\begin{proof}
The increment of the twist is
$(x_0+h_0)(y_0+h_0)^2-x_0y_0^2
=2h_0x_0y_0+h_0y_0^{2}+2h_0^{2}y_0+h_0^{2}x_0+h_0^{3}$.
Insert $\Delta_{(h,h)}f=\Delta_{(h,h)}g\cdot e_p(\text{increment})$ into the
box expression. The rank-one splitting of $\Delta_{(h,h)}g$
(Theorem~\ref{thm:D}) cancels pairwise against the box pattern:
$\alpha_h(x_1)\overline{\alpha_h(x_2)}\,\overline{\alpha_h(x_1)}\,
\alpha_h(x_2)=1$, likewise in $y$; the constant contributes
$e_p(\omega_h)\overline{e_p(\omega_h)}^{\,2}e_p(\omega_h)=1$; and every
single-variable phase dies, since $\phi(x_1)-\phi(x_2)-\phi(x_1)+\phi(x_2)=0$.
Only the mixed monomial survives, giving a bracket
$e_p\bigl(2h_0(u_0)(v_0)\bigr)$ with $u=x_1-x_2$, $v=y_1-y_2$ uniform, whose
average equals $\mathbb{P}(v_0=0)=p^{-1}$ whenever $2h_0\not\equiv0$. Of the
$p^{n}$ directions, $p^{n-1}$ have $h_0=0$ (value $1$) and $(p-1)p^{n-1}$
have $h_0\neq0$ (value $p^{-1}$); averaging yields $(2p-1)/p^{2}$. Since this
tends to $0$, the hypothesis fails at every fixed $\delta>0$ once $p$ is
large.
\end{proof}

Proposition~\ref{prop:collapse} refutes the cubic-twist route to a
counterexample: a family that violated the hypothesis could not contradict
the conjecture, and the averaged hypothesis is revealed as an instrument that
actively suppresses nondegenerate mixed-bilinear components of diagonal
increments. The companion strictness result, proved by the same four-point
difference technique, states that for $n\ge2$ the monomial phase
$e_p(\lambda x_1y_2)$, which enjoys pointwise rank-one diagonal increments
and hence lies in the rigidity family, cannot be written as
$a(x)b(y)c(x+y)$ with arbitrary unimodular one-variable factors: the assumed
identity forces, on the one hand,
$c(m+u+v)c(m)=e_p(\lambda u_1v_2)\,c(m+u)c(m+v)$ for all $m,u,v$, whose left
side is symmetric in $u\leftrightarrow v$ while its right side is not —
taking $\{u,v\}=\{e_1,e_2\}$ yields $e_p(\lambda)=1$, a contradiction. Thus
the phase-restricted target class is properly contained in the rigidity
family; under the unrestricted-quadratic reading of $(-1)^{q}$, by contrast,
the witness belongs to the transmitted class through $Q=\lambda x_1y_2$, and
the separation question re-enters through correlation, that is, through
Open Problem~(O1) below.

We now turn to the converse classification. Write points of $G\times G$ in
the coordinates $u=x+y$, $v=x-y$ and set
$\Theta(u,v)=f\bigl(\tfrac{u+v}{2},\tfrac{u-v}{2}\bigr)$, so that $f$ and
$\Theta$ determine each other. Under this pullback the pointwise rank-one
condition on diagonal increments becomes a multiplicative constraint on
$\Theta$ involving four evaluations arranged antipodally about their common
midpoint. An earlier version of this manuscript proposed the four-point
parallelogram equation \eqref{eq:parallelogram} below as a necessary
condition; that proposal is \emph{false}, and the failure is recorded
explicitly before stating the conjecture in its corrected, purely open form.

\begin{remark}[The parallelogram equation is not necessary]
Consider, for arbitrary $u,v,w,t\in G$, the equation
\begin{equation}\label{eq:parallelogram}
\Theta(u+w,\,v+t)\,\Theta(u-w,\,v-t)
\;=\;\Theta(u+t,\,v-w)\,\Theta(u-t,\,v+w).
\end{equation}
It is violated by members of the sufficiency family itself
(Theorem~\ref{thm:D}), so it is \emph{not} a necessary condition for
pointwise rank-one diagonal increments and cannot serve as the necessity
mechanism for the converse. Two-line counterexamples at $p=5$: for
$f(x,y)=e_p(x^2)$, i.e.\ $\Theta(u,v)=e_p\bigl((u+v)^2/4\bigr)$, at
$(u,v,w,t)=(1,-1,1,1)$ the left side of \eqref{eq:parallelogram} equals
$e_p(2)$ while the right side equals $1$; likewise for $f(x,y)=e_p(xy)$,
i.e.\ $\Theta(u,v)=e_p\bigl((u^2-v^2)/4\bigr)$, the two sides differ by the
nontrivial factor $e_p\bigl((w^2-t^2)/2\bigr)$ whenever $w^2\not\equiv t^2$.
Even the plain constituent $\Theta(u,v)=d(v)$ present in every family member
violates it, e.g.\ for $d(v)=e_p(v^2)$ whenever $t^2\not\equiv w^2$.
\end{remark}

\begin{conjecture}[Converse classification]\label{conj:Dp}
Suppose $f:G\times G\to\mathbb{C}$ is unimodular with pointwise rank-one
diagonal increments. Must $f$ belong to the closure of the family
$a(x)b(y)d(x-y)e_p(Q(x,y))$ with $Q$ quadratic? Sufficiency of the family
for pointwise rank-one increments is Theorem~\ref{thm:D}; the exact converse
classification is open, the defining rank-one increment law being the only
currently verified invariant, and no functional equation characterizing the
family — in particular not the parallelogram equation
\eqref{eq:parallelogram} displayed above, which its own members violate —
is presently known.
\end{conjecture}

\begin{corollary}[Misalignment between hypothesis and stated target]\label{cor:misalign}
The pointwise-hypothesis variant of the problem, which asks for structure
explaining rank-one diagonal increments, carries the natural target family
consisting of the functions $a(x)b(y)d(x-y)e_p(Q)$ with $Q$ quadratic,
whereas the transmitted statement names the functions of the form
$a(x)b(y)c(x+y)(-1)^{q}$. For $n\ge2$ these classes differ: the phase
$e_p(\lambda x_1y_2)$ belongs to the former family, taking $d\equiv1$ and
$Q=\lambda x_1y_2$, but lies outside the span of phases absorbable into
$a,b,c$ of the latter by Lemma~\ref{lem:noabsorb}. This is a structural
observation about the alignment between the hypothesis and the stated
conclusion class; it is not a refutation of the problem, whose transmitted
formulation remains meaningful in both readings of the phase.
\end{corollary}

\section{Numerical findings}\label{sec:numerics}

All numbers in this section were obtained by exact rational and modular
arithmetic, or by high-precision floating-point evaluation, with
deterministic pseudo-random initialization. The identity diagnostics for
Proposition~\ref{prop:A}, including the Parseval diagnostic that the total
Fourier mass of $B_f$ equals $1.0000000$ at all three sizes tested, were
already displayed in Table~\ref{tab:ident}.

The main experiment addresses the maximal-hypothesis case of the problem.
By Observation~\ref{obs:C} every unimodular diagonal function meets the
hypothesis with $\delta=1$, and by Lemma~\ref{lem:absorb} any quadratic
phase factor is redundant when $n=1$, so optimizing over arbitrary
unimodular $a,b,c$ already dominates the supremum over the full stated
class. The quantity computed is therefore
$M=\sup| \mathbb{E}_{x,y}\,d(x-y)\overline{a(x)}\,\overline{b(y)}\,
\overline{c(x+y)} |$ over unimodular $a,b,c$, for the target
$d(s)=e_p(s^{3})$ on $G=\mathbb{F}_p$ with $p\in\{5,7\}$, by alternating
phase maximization with one thousand random restarts and at most three
hundred sweeps per restart. The figures reported below are the largest
correlations found by local alternating maximization from these many random
restarts; a certified global optimum is neither claimed nor obtainable by
this method, so each value is a lower bound for the corresponding supremum
$M$. Results are collected in Table~\ref{tab:exp},
together with a control run in which $d$ was taken with uniformly random
exponents.

\begin{table}[ht]
\centering
\small
\begin{tabular}{lllr}
\toprule
prime & target $d$ & statistic & value \\
\midrule
$p=5$ & $e_5(s^{3})$ & largest found correlation  & $0.7236067977499792$ \\
$p=5$ & $e_5(s^{3})$ & mean over restarts & $0.7225012249409789$ \\
$p=5$ & $e_5(s^{3})$ & standard deviation & $0.01745$ \\
$p=7$ & $e_7(s^{3})$ & largest found correlation  & $0.677276973$ \\
$p=7$ & $e_7(s^{3})$ & mean over restarts & $0.675801400$ \\
$p=7$ & $e_7(s^{3})$ & standard deviation & $0.01399$ \\
\midrule
$p=5$ & random exponents & largest found correlation & $0.7236067977499792$ \\
$p=7$ & random exponents & largest found correlation & $0.816021079$ \\
\bottomrule
\end{tabular}
\caption{Largest correlations found by local alternating maximization against
unimodular factors at $n=1$: lower bounds for the suprema; a certified global
optimum is neither claimed nor obtainable by this method.}
\label{tab:exp}
\end{table}

As a sanity control, the quadratic target $d(s)=e_p(2s^{2})$ admits the
explicit witness $a=e_p(4x^{2})$, $b=e_p(4y^{2})$, $c=e_p(-2s^{2})$: since
$4x^{2}+4y^{2}-2(x+y)^{2}=2(x-y)^{2}$ identically over $\mathbb{F}_p$, the
integrand has constant modulus one and the optimal correlation equals
exactly $1$, attained by this witness.\footnote{A naive witness sometimes
quoted for this control, namely $a=e_p(2x^{2})$, $b=e_p(2y^{2})$,
$c=e_p(-s^{2})$, is wrong: it gives
$c\cdot d\cdot\overline{ab}=e_p(s^{2})$ after simplification, a pure Gauss
sum of magnitude $\sqrt{p}$, so the correlation evaluates to
$1/\sqrt{p}$ rather than to $1$. The corrected witness above restores the
exact value.}

The interpretation must be stated carefully. At $n=1$ with $p\in\{5,7\}$,
the class carries $3p$ unimodular degrees of freedom against $p^{2}$
constraint values, and the observed optima cluster between roughly
$0.68$ and $0.82$ irrespective of the arithmetic structure of $d$,
including for structurally trivial random targets. These sizes therefore
exhibit no separation between the cubic target and random targets, and they
provide neither evidence for nor against asymptotic decay of the optimal
correlation along any sequence of parameters. Whether there exists a
counterexample family, satisfying the hypothesis at $\delta=1$ while the
optimal correlation tends to zero along some sequence of primes or
dimensions, remains open.

\section{What remains open}\label{sec:open}

The following problems are stated precisely and remain open.

\begin{itemize}
\item[\textnormal{(O1)}] \emph{The twisted trilinear inequality.} Prove or
disprove the uniform lower bound formulated in Section~\ref{sec:diagonal}
for the display \eqref{eq:twistedcorr} over all unimodular $d$, with the
qualifications recorded there concerning the range of the twist for $n\ge2$
versus $n=1$.
\item[\textnormal{(O2)}] \emph{The converse classification.} Establish
Conjecture~\ref{conj:Dp}: whether pointwise rank-one diagonal increments
force membership in the family $a(x)b(y)d(x-y)e_p(Q)$.
\item[\textnormal{(O3)}] \emph{Decay and rank obstructions.} Determine
whether the supremal correlation decays along some sequence of parameters
$(p,n)$ for the cubic target $d=e_p(s^{3})$; relate the question to the
slice rank and partition rank of the kernel
$(s,m)\mapsto d(s)e_p(Q(s,m))$, in the spirit of the symmetric rank
formulation of cap-set bounds \cite{tao2016}, and to the true-complexity
analysis of systems of linear equations \cite{gowerswolf2010}.
\end{itemize}

It is instructive to record why naive approaches fail. Because the factors
$a,b,c$ are arbitrary bounded functions rather than phases of prescribed
polynomial form, the spectral reformulation of Proposition~\ref{prop:A},
which controls only aggregate Fourier mass, cannot by itself produce a
correlation with structured class members; and because the class is not
finite-dimensional, polynomial-method and rank-based arguments do not
parametrize its members. Any successful approach must exploit the
interplay between the diagonal-derivative structure of the hypothesis and
the three-fold additive structure of the target class, which is precisely
what the reductions of Sections~\ref{sec:calc} and~\ref{sec:rigidity}
isolate.

\begin{thebibliography}{9}
\raggedright

\bibitem{greentao2010}
B.~Green and T.~Tao,
An inverse theorem for the Gowers $U^{3}(G)$-norm,
Proc.\ Edinburgh Math.\ Soc.\ (2) \textbf{53} (2010), 81--132.

\bibitem{gowerswolf2010}
W.~T.~Gowers and J.~Wolf,
The true complexity of a system of linear equations,
Proc.\ London Math.\ Soc.\ (3) \textbf{100} (2010), 155--176.

\bibitem{tao2016}
T.~Tao,
A symmetric formulation of the Croot--Lev--Pach--Ellenberg--Gijswijt capset
bound,
electronic notes, 2016.

\bibitem{unsolvedmath}
UnsolvedMath problem archive, entry GREEN-028,
Gowers box norms over finite fields,
\url{https://www.unsolvedmath.com/problems/GREEN-028}.

\end{thebibliography}

\end{document}
